\documentclass[12pt, a4paper]{ctexart}
\usepackage{fontspec}
\usepackage{xcolor}
\usepackage{hyperref}
\usepackage{geometry}
\usepackage{enumitem}
\usepackage{parskip}
\usepackage{titlesec}
\usepackage{tcolorbox}
\tcbuselibrary{breakable}
\tcbset{autoparskip}

\usepackage{setspace}
\usepackage{listings}
\usepackage{amssymb}

\geometry{left=2.5cm, right=2.5cm, top=2.5cm, bottom=2.5cm}

\setmainfont{Times New Roman}
\setCJKmainfont{SimSun}

\hypersetup{
    colorlinks=true,
    linkcolor=blue!70!black,
    urlcolor=cyan,
    pdftitle={I/O-逐题精讲解义},
    pdfauthor={专升本-fifine老师}
}

\definecolor{myblue}{RGB}{30,100,200}
\definecolor{lightblue}{RGB}{230,240,255}
\definecolor{darkblue}{RGB}{0,50,120}

\newtcolorbox{bluebox}[1][]{
    breakable,
    colback=lightblue,
    colframe=myblue,
    coltitle=darkblue,
    fonttitle=\bfseries,
    title={#1},
    boxrule=0.8pt,
    arc=4pt,
    left=6pt,
    right=6pt,
    top=6pt,
    bottom=6pt,
    before skip=8pt,
    after skip=8pt
}

\usepackage{etoolbox}
\BeforeBeginEnvironment{bluebox}{\par\vfill}%
\AfterEndEnvironment{bluebox}{\par\vfill}%

\titleformat{\section}
  {\normalfont\Large\bfseries\color{myblue}}{\thesection}{1em}{}
\titlespacing*{\section}{0pt}{3.5ex plus 1ex minus .2ex}{1.5ex plus .2ex}

\title{\textcolor{myblue}{\textbf{I/O - 逐题精讲解义}}}
\author{专升本-fifine老师 教学组}
\date{2026年03月05日}

\lstset{
    language=Java,
    basicstyle=\ttfamily\small,
    keywordstyle=\color{blue},
    commentstyle=\color{green!60!black},
    stringstyle=\color{orange},
    numbers=left,
    numberstyle=\tiny\color{gray},
    frame=single,
    breaklines=true,
    columns=fullflexible,
    showstringspaces=false
}

\newcommand{\code}[1]{\texttt{#1}}

\begin{document}

\maketitle
\thispagestyle{empty}
\newpage

\section{I/O}

\begin{enumerate}[label=\textbf{\arabic*.}]

	\item \textbf{以下哪个类用于实现文件读取（字符流）？}
	      \begin{enumerate}[label=\Alph*.]
		      \item FileReader
		      \item FileWriter
		      \item BufferedReader
		      \item BufferedWriter
	      \end{enumerate}

	      \begin{bluebox}{答案与深度解析}
		      \textbf{答案：A. FileReader}

		      \textbf{【核心认知框架】}
		      \begin{itemize}[label={}, leftmargin=*, nosep]
			      \item \textbf{题目本质}：考查I/O流的分类和用途
			      \item \textbf{关键区分点}：字符流vs字节流、读vs写、基本流vs缓冲流
			      \item \textbf{命题人意图}：测试对FileReader基本用法的识别
		      \end{itemize}

		      \textbf{【选项原理深度拆解】}

		      \begin{itemize}[leftmargin=*, nosep, label={}]
			      \item[A] \textbf{FileReader — 正确（字符文件读取）}
			            \begin{itemize}[label={}, leftmargin=*, nosep]
				            \item \textbf{继承关系}：FileReader → InputStreamReader → Reader
				            \item \textbf{功能说明}：读取文本文件，按字符处理（自动处理编码）
				            \item \textbf{代码示例}：
				                  \begin{lstlisting}[language=Java]
import java.io.*;

public class FileReaderDemo {
    public static void main(String[] args) {
        try {
            // 创建FileReader
            FileReader reader = new FileReader("test.txt");

            // 方式1：逐字符读取
            int ch;
            while ((ch = reader.read()) != -1) {
                System.out.print((char) ch);  // int转char
            }
            reader.close();

            // 方式2：读取到字符数组
            char[] buffer = new char[1024];
            int len;
            FileReader reader2 = new FileReader("test.txt");
            while ((len = reader2.read(buffer)) != -1) {
                System.out.print(new String(buffer, 0, len));
            }
            reader2.close();
        } catch (IOException e) {
            e.printStackTrace();
        }
    }
}
// FileReader用于读取字符流文本文件
				      \end{lstlisting}
				            \item \textbf{字符编码}：使用平台默认编码，建议指定编码
				                  \begin{lstlisting}[language=Java]
// 推荐方式：指定编码
FileReader reader = new FileReader("test.txt",
    Charset.forName("UTF-8"));
// 或者
InputStreamReader isr = new InputStreamReader(
    new FileInputStream("test.txt"), "UTF-8");
				      \end{lstlisting}
			            \end{itemize}

			      \item[B] \textbf{FileWriter — 错误（字符文件写入）}
			            \begin{itemize}[label={}, leftmargin=*, nosep]
				            \item \textbf{继承关系}：FileWriter → OutputStreamWriter → Writer
				            \item \textbf{功能说明}：写入文本文件，按字符处理
				            \item \textbf{代码示例}：
				                  \begin{lstlisting}[language=Java]
FileWriter writer = new FileWriter("output.txt");
writer.write("Hello, World!");
writer.write("你好，世界；");
writer.close();

// 追加模式
FileWriter appendWriter = new FileWriter(
    "output.txt", true);  // true表示追加
appendWriter.write("追加的文本");
appendWriter.close();
// FileWriter用于写入字符流文本文件
				      \end{lstlisting}
			            \end{itemize}

			      \item[C] \textbf{BufferedReader — 错误（缓冲字符读取包装器）}
			            \begin{itemize}[label={}, leftmargin=*, nosep]
				            \item \textbf{功能说明}：包装Reader，提供缓冲功能，提高读取效率
				            \item \textbf{特别方法}：readLine()读取一行
				            \item \textbf{代码示例}：
				                  \begin{lstlisting}[language=Java]
// BufferedReader包装FileReader
BufferedReader br = new BufferedReader(
    new FileReader("test.txt"));

String line;
while ((line = br.readLine()) != null) {
    System.out.println(line);  // 逐行读取
}
br.close();

// 自动关闭资源的推荐写法
try (BufferedReader br = new BufferedReader(
        new FileReader("test.txt"))) {
    String line;
    while ((line = br.readLine()) != null) {
        System.out.println(line);
    }
}
// BufferedReader提供缓冲和按行读取功能
				      \end{lstlisting}
			            \end{itemize}

			      \item[D] \textbf{BufferedWriter — 错误（缓冲字符写入包装器）}
			            \begin{itemize}[label={}, leftmargin=*, nosep]
				            \item \textbf{功能说明}：包装Writer，提供缓冲功能，提高写入效率
				            \item \textbf{代码示例}：
				                  \begin{lstlisting}[language=Java]
BufferedWriter bw = new BufferedWriter(
    new FileWriter("output.txt"));

bw.write("第一行");
bw.newLine();  // 写入换行符
bw.write("第二行");
bw.write("第三行");
bw.flush();  // 刷新缓冲区
bw.close();
// BufferedWriter提供缓冲提高写入效率
				      \end{lstlisting}
			            \end{itemize}
		      \end{itemize}

		      \textbf{【I/O流分类】}
		      \begin{itemize}[label={}, nosep]
			      \item 按数据单位：字节流（InputStream/OutputStream）、字符流（Reader/Writer）
			      \item 按作用：输入流（读）、输出流（写）
			      \item 按功能：节点流（直接访问数据源）、处理流（包装节点流）
		      \end{itemize}

		      \textbf{【应背诵知识点】}
		      \begin{itemize}[label={}, nosep]
			      \item FileReader：读取字符文件（文本）
			      \item FileWriter：写入字符文件（文本）
			      \item BufferedReader：缓冲字符读取（可按行）
			      \item BufferedWriter：缓冲字符写入
			      \item 口诀：FileReader读字符文件，FileWriter写字符文件
		      \end{itemize}

		      \textbf{【命题趋势与实战策略】}
		      \textbf{考场秒杀}：看到"文件读取字符流"→FileReader
	      \end{bluebox}

\end{enumerate}

\end{document}
